\documentclass[aspectratio=169,11pt]{beamer}

\usetheme{metropolis}
\usepackage{booktabs}
\usepackage{graphicx}

\definecolor{navyblue}{RGB}{0,51,102}
\definecolor{oceanblue}{RGB}{0,102,179}
\definecolor{alertred}{RGB}{204,0,0}

\setbeamercolor{frametitle}{bg=navyblue, fg=white}
\setbeamercolor{progress bar}{fg=oceanblue}
\setbeamercolor{alerted text}{fg=alertred}

\setbeamertemplate{footline}{}
\setbeamertemplate{navigation symbols}{}

% ─────────────────────────────────────────────────────────────────────────────
\begin{document}

% ─────────────────────────────────────────────────────────────────────────────
\begin{frame}{The Problem: Vessels Go Dark and Disguise Their Activity}
  \begin{columns}[T]
    \begin{column}{0.52\textwidth}
      \textbf{90,000 vessels} of regulatory concern are active at any
      given moment — and the picture is incomplete by design.

      \bigskip
      \alert{75\% of industrial fishing vessels} intentionally disable
      their AIS transponders to avoid detection.

      \bigskip
      Even when a vessel \emph{is} visible, analysts must answer two
      questions from limited, noisy data:

      \medskip
      \begin{itemize}
        \item \textbf{What is it doing?}\\
              Fishing · Transiting · Loitering · Ship-to-Ship transfer
        \item \textbf{What type of vessel is it?}\\
              Trawler · Cargo · Naval · Unknown contact
      \end{itemize}

      \bigskip
      The data may be a full day of AIS pings — or a \textbf{single
      satellite image}.
    \end{column}
    \begin{column}{0.45\textwidth}
      \includegraphics[width=\textwidth]{../outputs/figures/B2_dark_period_uncertainty_cones.png}

      \smallskip
      \includegraphics[width=\textwidth]{../outputs/figures/C3_sts_rendezvous.png}
      {\small\textit{Top: position uncertainty after AIS loss.\\
      Bottom: ship-to-ship transfer — a key evasion tactic.}}
    \end{column}
  \end{columns}
\end{frame}

% ─────────────────────────────────────────────────────────────────────────────
\begin{frame}{The Scale: Where Fishing Actually Happens}
  \begin{columns}[T]
    \begin{column}{0.52\textwidth}
      The system is grounded in \textbf{real fleet data} from Global
      Fishing Watch — 2023 effort across four operationally critical
      regions.

      \bigskip
      \begin{tabular}{lrr}
        \toprule
        Region & Records & Fishing hrs \\
        \midrule
        Brazil EEZ         & 309,711 & 1,382,042 \\
        Philippine EEZ     & 61,460  & — \\
        Strait of Malacca  & 54,040  & — \\
        Gulf of Guinea     & 59,954  & — \\
        \bottomrule
      \end{tabular}

      \bigskip
      Fishing effort is \textbf{not uniformly distributed} — it clusters
      around productive grounds and shifts seasonally.

      \medskip
      Our synthetic simulator is calibrated against these real
      distributions, keeping speed profiles, gear-type ratios,
      and dark-event rates realistic.

      \medskip
      Speed profile KL-divergence vs.\ GFW real data:\\
      mean $D_{KL} < 0.15$ nats across all gear types.
    \end{column}
    \begin{column}{0.45\textwidth}
      \includegraphics[width=\textwidth]{../outputs/figures/B1_gfw_fishing_effort_asia_pacific.png}
      {\small\textit{GFW 2023 fishing effort — Asia-Pacific.
      Dense clusters mark productive grounds in the Philippine EEZ
      and South China Sea. The Strait of Malacca corridor is a
      key chokepoint for both fishing and transit traffic.}}
    \end{column}
  \end{columns}
\end{frame}

% ─────────────────────────────────────────────────────────────────────────────
\begin{frame}{The System: Classify Any Vessel from Any Observation}
  \begin{columns}[T]
    \begin{column}{0.52\textwidth}
      We built an \textbf{Activity Intelligence Engine} that works
      regardless of how much data is available.

      \bigskip
      \textbf{Inputs:} AIS, radar, satellite imagery (EO/SAR),
      or any combination — down to a single observation.

      \bigskip
      \textbf{Outputs per vessel:}
      \begin{itemize}
        \item Activity class (fishing / transit / anchored / STS / loitering)
        \item Vessel type (12 classes)
        \item Dark-vessel risk score (0--1)
        \item Calibrated confidence — conservative for sparse data
      \end{itemize}

      \bigskip
      \textbf{How it handles missing data:}\\
      A GPU-accelerated XGBoost model learns what each feature's
      \emph{absence} means — no imputation, no minimum track length
      required.

      \bigskip
      Trained and evaluated on \textbf{5,121 real vessels} from NOAA
      AIS data across three independent datasets.
    \end{column}
    \begin{column}{0.45\textwidth}
      \includegraphics[width=\textwidth]{../outputs/figures/B5_risk_score_dashboard.png}
      {\small\textit{Risk score dashboard — IUU and dark vessel scores
      per vessel. All flagged vessels are fishing-registered.}}

      \bigskip
      \includegraphics[width=\textwidth]{../outputs/brazil_eez/feature_importance.png}
      {\small\textit{Top features: speed mean, \% time slow, loiter index.}}
    \end{column}
  \end{columns}
\end{frame}

% ─────────────────────────────────────────────────────────────────────────────
\begin{frame}{Results: Robust Performance Across Regions and Sensors}
  \begin{columns}[T]
    \begin{column}{0.50\textwidth}
      \begin{tabular}{lc}
        \toprule
        \textbf{What we measure} & \textbf{Result} \\
        \midrule
        Activity classification $F_1$    & \textbf{0.900--1.000} \\
        Vessel type $F_1$ (12 classes)   & \textbf{0.969} \\
        Performance at $N=1$ observation & \textbf{same as $N=50$} \\
        \midrule
        Regions validated (synthetic)    & 4 \\
        \quad Brazil EEZ                 & $F_1 = 0.993$ \\
        \quad Strait of Malacca          & $F_1 = 0.999$ \\
        \quad Gulf of Guinea             & $F_1 = 0.992$ \\
        \quad Philippine EEZ             & $F_1 = 0.992$ \\
        \midrule
        Dark vessels flagged             & \textbf{49} \\
        \quad All fishing-registered     & \checkmark \\
        Training time (216k examples)   & \textbf{82 seconds} \\
        \bottomrule
      \end{tabular}

      \bigskip
      {\small The most significant open challenge is EO/SAR-only
      contacts where speed is unavailable.
      Geography-prior features and SAR foundation models are the
      next step.}
    \end{column}
    \begin{column}{0.47\textwidth}
      \includegraphics[width=\textwidth]{../outputs/figures/A2_partial_track_f1_vs_length.png}
      {\small\textit{Classification accuracy is flat from a single
      observation to a full 50-ping track — no minimum custody
      duration needed.}}

      \bigskip
      \includegraphics[width=\textwidth]{../outputs/brazil_eez/confusion_matrix.png}
      {\small\textit{Confusion matrix — Brazil EEZ.}}
    \end{column}
  \end{columns}
\end{frame}

\end{document}
